% Part: first-order-logic
% Chapter: completeness
% Section: downward-ls

\documentclass[../../../include/open-logic-section]{subfiles}

\begin{document}

\olfileid{fol}{com}{dls}
\olsection{De stelling van L\"owenheim--Skolem}

De stelling van L\"owenheim--Skolem zegt het volgende. Heeft een theorie
een oneindig model, dan heeft zij ook een model dat eindig of
!!{denumerable} is. Daaruit volgt meteen een grens van de
eerste-ordelogica: zij kan niet zeggen dat !!a{structure} een
!!{nonenumerable} grootte moet hebben. Neem namelijk !!a{sentence} of een
verzameling !!{sentence}s die in alle !!{nonenumerable} !!{structure}s
wordt vervuld. Die wordt ook in een !!{enumerable} structuur vervuld.

\begin{thm} 
\ollabel{thm:downward-ls} Als $\Gamma$ consistent is, heeft zij
!!a{enumerable} model. De verzameling is dus vervulbaar in !!a{structure}
waarvan het domein eindig of !!{denumerable} is.
\end{thm}

\begin{proof}
Stel dat $\Gamma$ consistent is. Het bewijs van de volledigheidsstelling
levert dan de !!{structure}~$\Struct M$. Haar domein $\Domain{M}$ is niet
groter dan de verzameling termen van de taal~$\Lang L$. De structuur
$\Struct M$ is dus hoogstens !!{denumerable}.
\end{proof}

\begin{thm}
\ollabel{noidentity-ls} Stel dat $\Gamma$ een consistente verzameling
!!{sentence}s is in de taal van de eerste-ordelogica zonder identiteit.
Dan heeft zij !!a{denumerable} model. Zij is dus vervulbaar in
!!a{structure} waarvan het domein oneindig en !!{enumerable} is.
\end{thm}

\begin{proof}
Stel dat $\Gamma$ consistent is en geen zinnen met identiteit bevat. Het
bewijs van de volledigheidsstelling levert dan de !!{structure}~$\Struct M$.
Haar domein $\Domain{M}$ is precies de verzameling termen van de
taal~$\Lang L'$. De structuur $\Struct{M}$ is dus !!{denumerable}, want
$\Trm[L']$ is dat ook.
\end{proof}

\begin{ex}[De paradox van Skolem]
De verzamelingenleer van Zermelo--Fraenkel~$\Log{ZFC}$ is een krachtig
raamwerk. We kunnen er vrijwel alle wiskundige uitspraken in doen, ook
over de grootte van verzamelingen. Zo bewijst $\Log{ZFC}$ dat de
verzameling~$\Real$ van reële getallen !!{nonenumerable} is. De theorie
bewijst ook de stelling van Cantor: de machtsverzameling van een
verzameling is groter dan die verzameling zelf. Als $\Log{ZFC}$
consistent is, zijn al haar modellen oneindig. Elk model bevat bovendien
!!{element}s waarvan de theorie zegt dat ze !!{nonenumerable} zijn. Een
voorbeeld is het element dat de stelling van~$\Log{ZFC}$ waar maakt dat
de machtsverzameling van de natuurlijke getallen bestaat. Toch heeft
$\Log{ZFC}$ volgens de stelling van L\"owenheim--Skolem ook
!!{enumerable} modellen. Zo'n model bevat ``!!{nonenumerable}''
verzamelingen, terwijl het model zelf !!{enumerable} is.
\end{ex}

\end{document}
